\documentclass[11pt,a4paper]{article}
\usepackage[top=1.7cm,bottom=1.5cm,left=1.5cm,right=1.5cm,headsep=8pt,footskip=16pt]{geometry}
\usepackage{enumitem}
\usepackage{amssymb}
\usepackage{array}
\usepackage[table]{xcolor}
\usepackage{fancyhdr}
\usepackage{lastpage}
\usepackage{multicol}
\usepackage{polyglossia}
\setdefaultlanguage[calendar=gregorian, numerals=maghrib]{arabic}
\setotherlanguage{english}
\IfFontExistsTF{Amiri}{%
  \newfontfamily\arabicfont[Script=Arabic, Scale=1.08]{Amiri}
}{%
  \newfontfamily\arabicfont[Script=Arabic, Scale=1.08, Path=../fonts/, Extension=.ttf,
    UprightFont=*-Regular, BoldFont=*-Bold, ItalicFont=*-Italic,
    BoldItalicFont=*-BoldItalic]{Amiri}
}
\IfFontExistsTF{AlBattar}{%
  \newfontfamily\khatt[Script=Arabic, Scale=1.15]{AlBattar}
}{%
  \newfontfamily\khatt[Script=Arabic, Scale=1.15, Path=../fonts/,
    BoldFont=ae_AlBattar.ttf]{ae_AlBattar.ttf}
}
\renewcommand{\headrulewidth}{0.8pt}
\renewcommand{\footrulewidth}{0.8pt}
\fancyhf{}
\rhead{\khatt المدرسة العليا للأساتذة ببوزريعة}
\lhead{\khatt ملحقة بومرداس}
\lfoot{\khatt الأستاذ: د. بن عبد الرزاق يوسف}
\cfoot{\thepage/\pageref{LastPage}}
\pagestyle{fancy}
\parindent 0pt
\setlength{\columnsep}{0.6cm}
\setlength{\columnseprule}{0.4pt}
\setlist[enumerate,1]{label=\textbf{\arabic*.}, leftmargin=*, itemsep=5pt, topsep=2pt}
\begin{document}

\linespread{1.08}\selectfont
\begin{center}
{\LARGE\khatt\textbf{الاختبار السادس عشر في الذكاء الاصطناعي}}\\[4pt]
{\normalsize وضعيات إشكالية ــ أسئلة متعددة الاختيارات ــ الفصول من 1 إلى 6}
\end{center}
\vspace{-2pt}\rule{\textwidth}{0.6pt}\\[2pt]
\textbf{التعليمة:} يتضمن الاختبار خمس وضعيات إشكالية. اقرأ كل وضعية بعناية، ثم اختر التحليل أو التصرف الأنسب
من بين الاختيارات الخمسة؛ لكل وضعية إجابة صحيحة واحدة فقط.
اكتب \textbf{حرف} الإجابة الصحيحة في جدول الإجابات في آخر الاختبار.
لكل وضعية 4 نقاط (المجموع 20 نقطة).
\vspace{4pt}

\noindent
{\large\khatt\textbf{الوضعية 1}}\hfill{\small(الفصل 1)}\par\nopagebreak\vspace{3pt}
\setlength{\fboxsep}{7pt}\colorbox{gray!10}{\parbox{\dimexpr\linewidth-14pt}{طلب أستاذ من أداة توليدية نصًا يصف مشاهد معركة بتفاصيل دموية، ليدرّس به خصائص شعر الحرب لتلاميذ السنة الرابعة متوسط، فاعتذرت الأداة عن التفاصيل العنيفة وعرضت بديلًا. اقترح زميل «حيلة» لتجاوز الرفض، ورأى آخر أن الأداة «جاهلة بالأدب». وهو يريد نصوصًا تخدم الدرس دون أن تصدم التلاميذ. ما الفهم الأدق لموقف الأداة؟}}\par\nopagebreak
\begin{itemize}[leftmargin=1.0cm, labelsep=6pt, itemsep=3pt, topsep=4pt, parsep=0pt]
  \item[\textbf{أ)}] الأداة تخضع لسياسات استخدام تقيّد بعض المحتوى، ورفضها لا يعني جهلها بالأدب؛ والأنسب توضيح السياق التعليمي وطلب تحليل نصوص شعرية أصيلة مناسبة للسن، لا التحايل على القيود.
  \item[\textbf{ب)}] يتخلى عن الأدوات التوليدية في تدريس الأدب نهائيًا، لأن أي أداة ترفض طلبًا واحدًا لا تصلح لأي مهمة أدبية، حتى تحليل النصوص الأصيلة أو اقتراح أسئلة الفهم حولها.
  \item[\textbf{ج)}] يستعمل حيلة الزميل بإيهام الأداة بأن الطلب لرواية خيالية، لأن الهدف التعليمي يبرّر أي وسيلة لتجاوز القيود، وتفاصيل العنف ضرورية لفهم شعر الحرب عند تلاميذ المتوسط.
  \item[\textbf{د)}] الرفض دليل على أن الأداة لا تعرف شعر الحرب أصلًا، لأن الأدوات ترفض فقط ما تجهله، فالحل استعمال أداة أخرى لا ترفض شيئًا، والاكتفاء بما تنتجه دون مراجعة مناسبته لسن التلاميذ.
  \item[\textbf{هـ)}] يكرر الطلب نفسه مرات كثيرة حتى تقبل الأداة، لأن الرفض قرار عشوائي يتغير بالإلحاح، وقبول الأداة في المرة العاشرة دليل على أن المحتوى مناسب للتلاميذ في كل الأحوال.
\end{itemize}
\par\vspace{6pt}
\noindent
{\large\khatt\textbf{الوضعية 2}}\hfill{\small(الفصل 2)}\par\nopagebreak\vspace{3pt}
\setlength{\fboxsep}{7pt}\colorbox{gray!10}{\parbox{\dimexpr\linewidth-14pt}{تتيح أداة يستعملها أستاذ إعدادًا يسمى «درجة الإبداع». لاحظ أنه عند رفعها تأتي اقتراحات العناوين متنوعة طريفة، لكن إجابات أسئلة القواعد تصبح أقل اتساقًا وتكثر فيها الأخطاء. ويريد ضبط الإعداد لكل مهمة من مهامه اليومية. ويرى أحد زملائه أن هذا الإعداد مجرد زخرفة لا أثر لها. ما التفسير والاستعمال الأنسب؟}}\par\nopagebreak
\begin{itemize}[leftmargin=1.0cm, labelsep=6pt, itemsep=3pt, topsep=4pt, parsep=0pt]
  \item[\textbf{أ)}] الإعداد يتحكم في سرعة الأداة، فرفعه يجعلها أسرع وأقل دقة، وخفضه يجعلها أبطأ وأدق، ولذلك يُخفض في كل المهام، لأن الدقة أهم من السرعة حتى في اقتراح العناوين.
  \item[\textbf{ب)}] الإعداد يتحكم في العشوائية، وخفضه إلى أدنى درجة يجعل إجابات القواعد صحيحة بالضرورة، فيكتفي بخفضه ويعتمد إجابات القواعد دون مراجعة، لأن الثبات دليل على الصحة.
  \item[\textbf{ج)}] الإعداد لا يؤثر إلا في طول الإجابة، والأخطاء في القواعد مصادفة، فيُضبط على أعلى درجة في كل المهام، لأن الإبداع مطلوب في القواعد كما هو مطلوب في اقتراح العناوين.
  \item[\textbf{د)}] الإعداد يتحكم في مدى عشوائية اختيار الكلمات أثناء التوليد؛ رفعه يزيد التنوع واحتمال الخطأ، وخفضه يجعل الإجابات أثبت؛ فيرفعه للعصف الذهني ويخفضه للقواعد، مع المراجعة في الحالتين.
  \item[\textbf{هـ)}] الإعداد يغيّر معارف النموذج نفسها، فرفعه يضيف إليه معلومات جديدة إبداعية، وخفضه يحذف بعض معارفه، ولذلك يُترك مرتفعًا دائمًا حتى يستفيد الأستاذ من كل معارف النموذج.
\end{itemize}
\par\vspace{6pt}
\noindent
{\large\khatt\textbf{الوضعية 3}}\hfill{\small(الفصل 3)}\par\nopagebreak\vspace{3pt}
\setlength{\fboxsep}{7pt}\colorbox{gray!10}{\parbox{\dimexpr\linewidth-14pt}{تريد أستاذة أسئلة فهم عن نص قرائي محدد في الكتاب المدرسي، فكتبت عنوان النص واسم كاتبه فقط، فجاءت أسئلة عن أحداث وشخصيات غير موجودة في النص. وتحتاج الأسئلة لحصة الغد. ويرى زميلها أن الأداة «تعرف كل نصوص الكتاب». والزميل نفسه جرّب ذلك مع نص آخر ولم يتحقق من النتيجة. أي طلب يعالج المشكلة؟}}\par\nopagebreak
\begin{itemize}[leftmargin=1.0cm, labelsep=6pt, itemsep=3pt, topsep=4pt, parsep=0pt]
  \item[\textbf{أ)}] تنسخ النص في الطلب، لكنها تطلب أسئلة «عن النص وعن كل ما يتصل به من معارف»، ثم تعتمد الأسئلة دون مراجعة، لأن وجود النص في الطلب يمنع الأداة من إضافة أي عنصر غير موجود فيه.
  \item[\textbf{ب)}] تنسخ النص كاملًا في الطلب، وتطلب أسئلة تعتمد عليه وحده مع ذكر الفقرة التي تُستخرج منها كل إجابة، ثم تتحقق من كل سؤال بالرجوع إلى النص، لأن الأداة لا تملك النص المدرسي بالضرورة.
  \item[\textbf{ج)}] تضيف إلى العنوان واسم الكاتب رقم الصفحة واسم الكتاب المدرسي، لأن هذه البيانات تجعل الأداة تسترجع النص حرفيًا من ذاكرتها، فتأتي الأسئلة مطابقة له دون حاجة إلى نسخه في الطلب.
  \item[\textbf{د)}] تطلب من الأداة كتابة نص جديد بالعنوان نفسه ثم أسئلة عنه، وتستعمل النص الجديد بدل نص الكتاب، لأن النص المولَّد يناسب الأسئلة، والتلاميذ لا يلاحظون الفرق بين النصين.
  \item[\textbf{هـ)}] تكتب «لا تخترع أحداثًا» مع العنوان واسم الكاتب فقط، لأن هذا النهي يكفي لتلتزم الأداة بالنص الأصلي الذي تعرفه، فتحذف الأحداث المختلقة وتُبقي الأسئلة المطابقة للكتاب.
\end{itemize}
\par\vspace{6pt}
\noindent
{\large\khatt\textbf{الوضعية 4}}\hfill{\small(الفصل 4)}\par\nopagebreak\vspace{3pt}
\setlength{\fboxsep}{7pt}\colorbox{gray!10}{\parbox{\dimexpr\linewidth-14pt}{تمنع أستاذة تلاميذها من استعمال الذكاء الاصطناعي في واجباتهم، بينما تُعدّ هي أغلب تمارينها ونصوصها بهذه الأدوات دون أن تخبرهم. اكتشف التلاميذ ذلك، فاحتجّوا بأن الأمر «ازدواجية». وتريد الأستاذة موقفًا متسقًا أخلاقيًا وتربويًا. وتخشى أن تفقد ثقة تلاميذها، وتريد أن تشرح لهم موقفها بوضوح وصدق في الحصة القادمة. ما الموقف الأنسب؟}}\par\nopagebreak
\begin{itemize}[leftmargin=1.0cm, labelsep=6pt, itemsep=3pt, topsep=4pt, parsep=0pt]
  \item[\textbf{أ)}] لا ازدواجية، لأن الأستاذ غير مطالب بالإفصاح لتلاميذه عن أدوات عمله، والتلاميذ لا يحق لهم مساءلته، فتستمر في المنع العام دون تعليل وفي الاستعمال دون إفصاح.
  \item[\textbf{ب)}] تُفصح عن استعمالها فقط، وتُبقي المنع العام دون تعليل أو تحديد، لأن الإفصاح وحده يرفع الازدواجية، ولا حاجة إلى شرح الفرق بين أهداف عملها وأهداف واجبات التلاميذ، فالشفافية تتحقق بمجرد الاعتراف بالاستعمال.
  \item[\textbf{ج)}] التلاميذ محقون تمامًا، فتسمح لهم باستعمال الأدوات في كل الواجبات دون ضوابط، لأن استعمال الأستاذة لها يعني أن الواجبات لم تعد تقيس شيئًا يحتاج التلميذ إلى بنائه بنفسه، فيكتسبون بذلك الحرية والمسؤولية معًا.
  \item[\textbf{د)}] تتوقف هي عن استعمال الأدوات نهائيًا وتُبقي المنع، لأن الاتساق الأخلاقي يتحقق بالامتناع التام من الجميع، والإفصاح وشرح الضوابط لا يكفيان لمعالجة الإشكال.
  \item[\textbf{هـ)}] الفرق مشروع إذا كان هدف الواجب بناء كفاءة التلميذ نفسه، لكن الشفافية تقتضي أن تُفصح عن استعمالها وتشرح ضوابطه، وأن تحدد للتلاميذ متى يُسمح بالأداة وكيف، بدل منع عام غير معلّل.
\end{itemize}
\par\vspace{6pt}
\noindent
{\large\khatt\textbf{الوضعية 5}}\hfill{\small(الفصل 6)}\par\nopagebreak\vspace{3pt}
\setlength{\fboxsep}{7pt}\colorbox{gray!10}{\parbox{\dimexpr\linewidth-14pt}{استعمل أستاذ أداة لتلوين صورة قديمة بالأبيض والأسود من أحداث 8 ماي 1945 و«تحسين دقتها» لعرضها في درس. جاءت الصورة الملوّنة واضحة مؤثرة، لكنه لاحظ أن الأداة أضافت تفاصيل في الوجوه واللافتات لم تكن في الصورة الأصلية. ويريد أن يستفيد من الصورة دون أن يضلل تلاميذه حول الوثيقة التاريخية. ما التصرف الأنسب؟}}\par\nopagebreak
\begin{itemize}[leftmargin=1.0cm, labelsep=6pt, itemsep=3pt, topsep=4pt, parsep=0pt]
  \item[\textbf{أ)}] يعرض النسخة الملونة مع تنبيه عام إلى أنها «ملوّنة»، دون الإشارة إلى التفاصيل المضافة، لأن التلوين هو التعديل الوحيد المهم، والإضافات في الوجوه واللافتات طفيفة لا تغيّر الوثيقة، ولا تستحق تنبيهًا خاصًا.
  \item[\textbf{ب)}] يعرض النسخة الملونة ويحذف الأصل من عرضه، لأن الصورة القديمة رديئة لا فائدة منها، والأداة استرجعت التفاصيل الحقيقية التي كانت موجودة في المشهد عند التصوير، فيرى التلاميذ الحدث كما وقع فعلًا.
  \item[\textbf{ج)}] يعرض الأصل بوصفه الوثيقة، ويقدّم النسخة المعالَجة إن استعملها على أنها تصوّر تقريبي، مع التنبيه إلى أن الأداة أضافت تفاصيل متخيَّلة، لأن «التحسين» الآلي قد يختلق ما ليس في الوثيقة.
  \item[\textbf{د)}] يتخلى عن عرض أي صورة تاريخية، لأن كل صورة قديمة معرّضة للتحريف، ولا يمكن التمييز بين الأصل والنسخة المعالَجة، حتى عند عرضهما معًا مع التنبيه.
  \item[\textbf{هـ)}] يعرض النسخة الملونة وحدها على أنها الصورة الأصلية، لأن التلوين والتحسين يقرّبان الحدث من التلاميذ، والتفاصيل المضافة تكمل ما أخفاه ضعف جودة الصورة القديمة.
\end{itemize}
\par\vspace{6pt}
\begin{center}
{\khatt\textbf{جدول الإجابات}}\\[3pt]
\renewcommand{\arraystretch}{1.6}
\begin{tabular}{|>{\centering\arraybackslash}m{1.9cm}|*{5}{>{\centering\arraybackslash}m{1.2cm}|}}
\hline
\textbf{رقم السؤال} & \textbf{1} & \textbf{2} & \textbf{3} & \textbf{4} & \textbf{5} \\ \hline
\textbf{الإجابة} &  &  &  &  &  \\[8pt] \hline
\end{tabular}

\vspace{8pt}
\renewcommand{\arraystretch}{2.2}
\begin{tabular}{|>{\centering\arraybackslash}m{7.5cm}|>{\centering\arraybackslash}m{3.2cm}|>{\centering\arraybackslash}m{3.2cm}|}
\hline
\textbf{الاسم واللقب} & \textbf{الفوج} & \textbf{العلامة} \\ \hline
 & & \hfill /20 \\ \hline
\end{tabular}
\end{center}
\end{document}
